% --- Archon physics formalization source begin ---
% archon:physics
% archon:covers PhyXMiniProblems/problem_phyx_mini_0678.lean
% archon:source-report reports/phyx_mini/problem_phyx_mini_0678.source.json

\chapter{Physics problem phyx\_mini\_0678}
\label{ch:PhyXMiniProblems_problem_phyx_mini_0678}

\paragraph{Problem source.}
\#\# Physical scenario

A snow-covered ski slope makes an angle of \textbackslash{}( 35.0\textasciicircum{}\textbackslash{}circ \textbackslash{}) with the horizontal. When a ski jumper plummets onto the hill, a parcel of splashed snow is thrown up to a maximum displacement of \textbackslash{}( 1.50 \textbackslash{}text\{ m\} \textbackslash{}) at \textbackslash{}( 16.0\textasciicircum{}\textbackslash{}circ \textbackslash{}) from the vertical in the uphill direction as shown in figure.

\#\# Question

Find the components of its maximum displacement perpendicular to the surface.

\#\# Answer choices

- A: A: \textbackslash{}( 0.178 \textbackslash{}text\{ m\} \textbackslash{})

- B: B: \textbackslash{}( 0.851 \textbackslash{}text\{ m\} \textbackslash{})

- C: C: \textbackslash{}( 0.944 \textbackslash{}text\{ m\} \textbackslash{})

- D: D: \textbackslash{}( 0.781 \textbackslash{}text\{ m\} \textbackslash{})

\#\# Figure caption (auxiliary; use the image as primary evidence)

The image depicts a skier going down a sloped incline. The skier is wearing a helmet and holding ski poles. The slope is angled at 35.0° from the horizontal. There is a black arrow pointing upwards and away from the slope, representing a trajectory or force, and it forms an angle of 16.0° with a vertical dashed line. The overall scene illustrates a concept from physics, likely related to motion on an inclined plane and projectile motion.

\#\# Dataset metadata

Kinematics; Spatial Relation Reasoning, Physical Model Grounding Reasoning

\paragraph{Recorded answer/context.}
C: C: \textbackslash{}( 0.944 \textbackslash{}text\{ m\} \textbackslash{})

\paragraph{Figure/image path.}
/root/proposal\_for\_physic/science-mango/phyx\_mini\_run/phyx\_data/test\_image/678.png

\paragraph{Formalization target.}
create a compiling Lean file with sorry bodies at `PhyXMiniProblems/problem\_phyx\_mini\_0678.lean`. The Lean declarations must preserve the physical quantities, dimensions or dimensional roles, figure labels, governing-law hypotheses, and final relation expressed by this problem.
Use Mathlib/Physlib names found through LeanExplore where available. If a physics API is missing, introduce faithful local abstractions rather than scalar placeholder aliases.

\begin{theorem}[Physics formalization target]
\label{thm:physics:phyx_mini_0678:target}
\lean{PhyXMiniProblems.ProblemPhyXMini0678.problem_phyx_mini_0678}
\uses{def:physics:phyx-mini-0678:phyxminiproblems-problemphyxmini0678-lengthinmeters, def:physics:phyx-mini-0678:phyxminiproblems-problemphyxmini0678-degreestoradians, def:physics:phyx-mini-0678:phyxminiproblems-problemphyxmini0678-skislopedisplacementsetup, def:physics:phyx-mini-0678:phyxminiproblems-problemphyxmini0678-matchesproblemandprimaryfigure, def:physics:phyx-mini-0678:phyxminiproblems-problemphyxmini0678-obeysperpendicularprojectionlaw, def:physics:phyx-mini-0678:phyxminiproblems-problemphyxmini0678-recordedanswerchoice, def:physics:phyx-mini-0678:phyxminiproblems-problemphyxmini0678-matchesdisplayedanswertothreedecimalplaces}
The assigned autoformalize agent should translate this physics problem into Lean declarations in the covered file, with theorem and lemma proof bodies written as `by sorry`.
\end{theorem}
\begin{proof}
This is an autoformalization task, not a proof task. Produce statements that can later be proved by the physics prover without weakening the physical model.
\end{proof}

% --- Archon declaration topology begin ---
\section{Declaration topology}
\begin{definition}[Diagram Plane]
  \label{def:physics:phyx-mini-0678:phyxminiproblems-problemphyxmini0678-diagramplane}
  \lean{PhyXMiniProblems.ProblemPhyXMini0678.DiagramPlane}
  The two-dimensional Euclidean plane containing the slope and snow arrow
\end{definition}
\begin{proof}
  This is the declared model definition of Diagram Plane.
\end{proof}
\begin{definition}[Planar Displacement Quantity]
  \label{def:physics:phyx-mini-0678:phyxminiproblems-problemphyxmini0678-planardisplacementquantity}
  \lean{PhyXMiniProblems.ProblemPhyXMini0678.PlanarDisplacementQuantity}
  \uses{def:physics:phyx-mini-0678:phyxminiproblems-problemphyxmini0678-diagramplane}
  A unit-independent planar displacement carrying physical length dimension
\end{definition}
\begin{proof}
  This is the declared model definition of Planar Displacement Quantity.
\end{proof}
\begin{definition}[Scalar Length Quantity]
  \label{def:physics:phyx-mini-0678:phyxminiproblems-problemphyxmini0678-scalarlengthquantity}
  \lean{PhyXMiniProblems.ProblemPhyXMini0678.ScalarLengthQuantity}
  A signed unit-independent scalar component carrying length dimension
\end{definition}
\begin{proof}
  This is the declared model definition of Scalar Length Quantity.
\end{proof}
\begin{definition}[displacement Vector In Meters]
  \label{def:physics:phyx-mini-0678:phyxminiproblems-problemphyxmini0678-displacementvectorinmeters}
  \lean{PhyXMiniProblems.ProblemPhyXMini0678.displacementVectorInMeters}
  \uses{def:physics:phyx-mini-0678:phyxminiproblems-problemphyxmini0678-diagramplane, def:physics:phyx-mini-0678:phyxminiproblems-problemphyxmini0678-planardisplacementquantity}
  Read a planar displacement in coherent SI units, i.e. metres
\end{definition}
\begin{proof}
  This is the declared model definition of displacement Vector In Meters.
\end{proof}
\begin{definition}[length In Meters]
  \label{def:physics:phyx-mini-0678:phyxminiproblems-problemphyxmini0678-lengthinmeters}
  \lean{PhyXMiniProblems.ProblemPhyXMini0678.lengthInMeters}
  \uses{def:physics:phyx-mini-0678:phyxminiproblems-problemphyxmini0678-scalarlengthquantity}
  Read a signed scalar length in coherent SI units, i.e. metres
\end{definition}
\begin{proof}
  This is the declared model definition of length In Meters.
\end{proof}
\begin{definition}[displacement Magnitude In Meters]
  \label{def:physics:phyx-mini-0678:phyxminiproblems-problemphyxmini0678-displacementmagnitudeinmeters}
  \lean{PhyXMiniProblems.ProblemPhyXMini0678.displacementMagnitudeInMeters}
  \uses{def:physics:phyx-mini-0678:phyxminiproblems-problemphyxmini0678-planardisplacementquantity, def:physics:phyx-mini-0678:phyxminiproblems-problemphyxmini0678-displacementvectorinmeters}
  The Euclidean magnitude of a planar displacement's metre readout
\end{definition}
\begin{proof}
  This is the declared model definition of displacement Magnitude In Meters.
\end{proof}
\begin{definition}[degrees To Radians]
  \label{def:physics:phyx-mini-0678:phyxminiproblems-problemphyxmini0678-degreestoradians}
  \lean{PhyXMiniProblems.ProblemPhyXMini0678.degreesToRadians}
  Convert a dimensionless angle readout in degrees to radians
\end{definition}
\begin{proof}
  This is the declared model definition of degrees To Radians.
\end{proof}
\begin{definition}[unit Vector At Angle Degrees]
  \label{def:physics:phyx-mini-0678:phyxminiproblems-problemphyxmini0678-unitvectoratangledegrees}
  \lean{PhyXMiniProblems.ProblemPhyXMini0678.unitVectorAtAngleDegrees}
  \uses{def:physics:phyx-mini-0678:phyxminiproblems-problemphyxmini0678-diagramplane, def:physics:phyx-mini-0678:phyxminiproblems-problemphyxmini0678-degreestoradians}
  The unit direction at a counterclockwise angle from the positive horizontal axis. Coordinate 0 is horizontal and coordinate 1 is vertical
\end{definition}
\begin{proof}
  This is the declared model definition of unit Vector At Angle Degrees.
\end{proof}
\begin{definition}[Figure Feature]
  \label{def:physics:phyx-mini-0678:phyxminiproblems-problemphyxmini0678-figurefeature}
  \lean{PhyXMiniProblems.ProblemPhyXMini0678.FigureFeature}
  Named visual features in image phyx\_data/test\_image/678.png
\end{definition}
\begin{proof}
  This is the declared model definition of Figure Feature.
\end{proof}
\begin{definition}[Figure Direction]
  \label{def:physics:phyx-mini-0678:phyxminiproblems-problemphyxmini0678-figuredirection}
  \lean{PhyXMiniProblems.ProblemPhyXMini0678.FigureDirection}
  Qualitative directions distinguished by the problem and its figure
\end{definition}
\begin{proof}
  This is the declared model definition of Figure Direction.
\end{proof}
\begin{definition}[Ski Slope Figure]
  \label{def:physics:phyx-mini-0678:phyxminiproblems-problemphyxmini0678-skislopefigure}
  \lean{PhyXMiniProblems.ProblemPhyXMini0678.SkiSlopeFigure}
  \uses{def:physics:phyx-mini-0678:phyxminiproblems-problemphyxmini0678-figurefeature, def:physics:phyx-mini-0678:phyxminiproblems-problemphyxmini0678-figuredirection}
  Literal qualitative and numerical content of the supplied raster. The angle labels are dimensionless degree readouts; the physical displacement is stored separately in SkiSlopeDisplacementSetup
\end{definition}
\begin{proof}
  This is the declared model definition of Ski Slope Figure.
\end{proof}
\begin{definition}[Ski Slope Displacement Setup]
  \label{def:physics:phyx-mini-0678:phyxminiproblems-problemphyxmini0678-skislopedisplacementsetup}
  \lean{PhyXMiniProblems.ProblemPhyXMini0678.SkiSlopeDisplacementSetup}
  \uses{def:physics:phyx-mini-0678:phyxminiproblems-problemphyxmini0678-diagramplane, def:physics:phyx-mini-0678:phyxminiproblems-problemphyxmini0678-planardisplacementquantity, def:physics:phyx-mini-0678:phyxminiproblems-problemphyxmini0678-scalarlengthquantity, def:physics:phyx-mini-0678:phyxminiproblems-problemphyxmini0678-skislopefigure}
  Independent physical and geometrical quantities for the problem. The scalar perpendicularComponent is an observable to be determined by the projection law; no numerical answer is assigned to it here
\end{definition}
\begin{proof}
  This is the declared model definition of Ski Slope Displacement Setup.
\end{proof}
\begin{definition}[Matches Problem And Primary Figure]
  \label{def:physics:phyx-mini-0678:phyxminiproblems-problemphyxmini0678-matchesproblemandprimaryfigure}
  \lean{PhyXMiniProblems.ProblemPhyXMini0678.MatchesProblemAndPrimaryFigure}
  \uses{def:physics:phyx-mini-0678:phyxminiproblems-problemphyxmini0678-displacementvectorinmeters, def:physics:phyx-mini-0678:phyxminiproblems-problemphyxmini0678-displacementmagnitudeinmeters, def:physics:phyx-mini-0678:phyxminiproblems-problemphyxmini0678-unitvectoratangledegrees, def:physics:phyx-mini-0678:phyxminiproblems-problemphyxmini0678-skislopedisplacementsetup}
  Problem and primary-image readouts. Since the slope descends to the right, its uphill tangent has polar angle 180° - 35°, while its outward normal has polar angle 90° - 35°. The snow arrow lies left (uphill) of vertical and therefore has polar angle 90° + 16°. This structure fixes only the stated displacement and figure geometry. It does not mention the requested perpendicular component or any answer choice
\end{definition}
\begin{proof}
  This is the declared model definition of Matches Problem And Primary Figure.
\end{proof}
\begin{definition}[Obeys Perpendicular Projection Law]
  \label{def:physics:phyx-mini-0678:phyxminiproblems-problemphyxmini0678-obeysperpendicularprojectionlaw}
  \lean{PhyXMiniProblems.ProblemPhyXMini0678.ObeysPerpendicularProjectionLaw}
  \uses{def:physics:phyx-mini-0678:phyxminiproblems-problemphyxmini0678-displacementvectorinmeters, def:physics:phyx-mini-0678:phyxminiproblems-problemphyxmini0678-lengthinmeters, def:physics:phyx-mini-0678:phyxminiproblems-problemphyxmini0678-skislopedisplacementsetup}
  The standard Euclidean law for resolving a displacement perpendicular to a surface. The selected outward normal and uphill tangent form an orthonormal frame, and the signed normal component is their inner-product projection. This is a general governing relation and contains no requested numerical value
\end{definition}
\begin{proof}
  This is the declared model definition of Obeys Perpendicular Projection Law.
\end{proof}
\begin{lemma}[perpendicular Component exact]
  \label{lem:physics:phyx-mini-0678:phyxminiproblems-problemphyxmini0678-perpendicularcomponent-exact}
  \lean{PhyXMiniProblems.ProblemPhyXMini0678.perpendicularComponent_exact}
  \uses{def:physics:phyx-mini-0678:phyxminiproblems-problemphyxmini0678-lengthinmeters, def:physics:phyx-mini-0678:phyxminiproblems-problemphyxmini0678-degreestoradians, def:physics:phyx-mini-0678:phyxminiproblems-problemphyxmini0678-skislopedisplacementsetup, def:physics:phyx-mini-0678:phyxminiproblems-problemphyxmini0678-matchesproblemandprimaryfigure, def:physics:phyx-mini-0678:phyxminiproblems-problemphyxmini0678-obeysperpendicularprojectionlaw}
  The angle between the arrow and the outward surface normal is 35° + 16° = 51°, so the exact perpendicular component is 1.50 cos 51° metres
\end{lemma}
\begin{proof}
  The conclusion follows from the governing relations and figure readouts named in the statement of perpendicular Component exact.
\end{proof}
\begin{definition}[Answer Choice]
  \label{def:physics:phyx-mini-0678:phyxminiproblems-problemphyxmini0678-answerchoice}
  \lean{PhyXMiniProblems.ProblemPhyXMini0678.AnswerChoice}
  Labels of the four answers printed with the problem
\end{definition}
\begin{proof}
  This is the declared model definition of Answer Choice.
\end{proof}
\begin{definition}[displayed Component In Meters]
  \label{def:physics:phyx-mini-0678:phyxminiproblems-problemphyxmini0678-displayedcomponentinmeters}
  \lean{PhyXMiniProblems.ProblemPhyXMini0678.displayedComponentInMeters}
  \uses{def:physics:phyx-mini-0678:phyxminiproblems-problemphyxmini0678-answerchoice}
  The displacement component, in metres, printed beside each answer label
\end{definition}
\begin{proof}
  This is the declared model definition of displayed Component In Meters.
\end{proof}
\begin{definition}[recorded Answer Choice]
  \label{def:physics:phyx-mini-0678:phyxminiproblems-problemphyxmini0678-recordedanswerchoice}
  \lean{PhyXMiniProblems.ProblemPhyXMini0678.recordedAnswerChoice}
  \uses{def:physics:phyx-mini-0678:phyxminiproblems-problemphyxmini0678-answerchoice}
  The answer label recorded in the supplied dataset metadata
\end{definition}
\begin{proof}
  This is the declared model definition of recorded Answer Choice.
\end{proof}
\begin{definition}[Matches Displayed Answer To Three Decimal Places]
  \label{def:physics:phyx-mini-0678:phyxminiproblems-problemphyxmini0678-matchesdisplayedanswertothreedecimalplaces}
  \lean{PhyXMiniProblems.ProblemPhyXMini0678.MatchesDisplayedAnswerToThreeDecimalPlaces}
  \uses{def:physics:phyx-mini-0678:phyxminiproblems-problemphyxmini0678-lengthinmeters, def:physics:phyx-mini-0678:phyxminiproblems-problemphyxmini0678-skislopedisplacementsetup, def:physics:phyx-mini-0678:phyxminiproblems-problemphyxmini0678-answerchoice, def:physics:phyx-mini-0678:phyxminiproblems-problemphyxmini0678-displayedcomponentinmeters}
  Agreement with a value displayed to three decimal places in metres. The half-unit in the last shown digit is 0.0005 m, avoiding the false assertion that the rounded decimal is exactly the trigonometric value
\end{definition}
\begin{proof}
  This is the declared model definition of Matches Displayed Answer To Three Decimal Places.
\end{proof}
% --- Archon declaration topology end ---
% --- Archon physics formalization source end ---
